% Research manuscript developed from the supplied Springer Nature sn-article.tex.
% The official Springer Nature sn-jnl.cls is required for compilation.
% This is a working draft. No experimental results have been invented.
% Literature and standards: see self-contained thebibliography section.
\documentclass[pdflatex,sn-mathphys-num]{sn-jnl}

\usepackage{graphicx}
\usepackage{amsmath,amssymb,amsfonts}
\usepackage{booktabs}
\usepackage{multirow}
\usepackage{algorithm}
\usepackage{algpseudocode}
\usepackage{url}
\usepackage{array}

\newcommand{\vect}[1]{\boldsymbol{#1}}
\newcommand{\mat}[1]{\mathbf{#1}}
\newcommand{\norm}[1]{\left\lVert #1\right\rVert}
\newcommand{\trans}{^{\mathsf T}}
\newcommand{\prob}{\mathbb{P}}
\newcommand{\pending}{\emph{Pending experimental evaluation}}
\newcommand{\todo}[1]{\textbf{[TO BE COMPLETED: #1]}}

\begin{document}

\title[Integrity Assessment of Galileo and 5G DL-TDOA]{A Comparative Integrity Assessment Framework for Galileo and 5G DL-TDOA Positioning in Urban Canyons}

% Authorship, institutional affiliations, and corresponding email must be
% confirmed by the authors before submission. No co-authors are invented.
\author*[1]{\fnm{David} \sur{S\'anchez-Heredero}}
\affil*[1]{\orgname{Affiliation to be confirmed}}

\abstract{Global Navigation Satellite System (GNSS) and fifth-generation (5G)
New Radio (NR) positioning offer complementary geometries in urban environments,
but positioning accuracy alone does not quantify the trustworthiness of a location
estimate. This paper develops a common integrity-assessment methodology for
single-frequency Galileo pseudorange positioning and 5G NR downlink time
difference of arrival (DL-TDOA). For Galileo, a linearized least-squares model
provides the nominal position covariance, while residual-based receiver
autonomous integrity monitoring (RAIM) characterizes single-measurement fault
detectability through central and noncentral chi-square distributions. For 5G,
a nonlinear weighted least-squares estimator accounts for the non-diagonal
covariance induced by common-reference TDOA measurements. Conditional horizontal
and vertical protection bounds are derived for Gaussian errors and explicitly
bounded timing biases. Both approaches are designed for epoch-by-epoch assessment
within a shared simulated urban-canyon trajectory. The distinction between
model-conditional bounds, fault-detection capability, and a fully allocated
integrity-risk guarantee is maintained throughout. Experimental evaluation,
including protection-level coverage, availability, fault sensitivity, and
sensitivity to error-model assumptions, is specified but remains to be completed.
An extension to hybrid Galileo--5G extended Kalman filtering is outlined as
future work.}

\keywords{Galileo, 5G NR, DL-TDOA, positioning integrity, protection level,
RAIM, weighted least squares, urban canyon, sensor fusion}

\maketitle

\section{Introduction}\label{sec:intro}

Reliable position estimation in urban environments requires more than low
root-mean-square error. Satellite signals can be attenuated or blocked by
buildings, whereas terrestrial 5G signals offer alternative transmitter
geometries but are themselves vulnerable to unfavorable geometry,
non-line-of-sight (NLOS) propagation, and imperfect network synchronization.
An estimator may therefore report a plausible location while its actual
position error is substantially larger than predicted.

Positioning \emph{accuracy} characterizes the magnitude of an estimation
error, whereas \emph{integrity} concerns the probability that a positioning
system provides misleading information without an appropriate warning.
Protection levels are meaningful only when linked to stated statistical
assumptions, fault hypotheses, and risk allocation. Classical receiver
autonomous integrity monitoring (RAIM) uses redundant GNSS measurements
to test measurement consistency \cite{Brown1992,WalterEnge1995}. More recently,
3GPP has introduced positioning integrity support for radio-access-technology
(RAT)-dependent positioning \cite{ETSI38305}, and dedicated frameworks for
5G positioning integrity have been investigated \cite{Bartoletti2025}.

The starting point of this work is an existing simulation of a single-frequency
Galileo receiver, a PRS-based 5G DL-TDOA positioning system, and a hybrid
extended Kalman filter (EKF) in an urban corridor in Madrid
\cite{SanchezHerederoThesis}. That study focused primarily on positioning
accuracy and availability. The present manuscript adds an explicit integrity
analysis while retaining the original scenario and differentiating existing
simulation components from the additional integrity methodology.

The methodological contributions investigated here are: (i) a directly
comparable local East--North--Up (ENU) uncertainty representation for Galileo
and 5G; (ii) a residual-based single-fault detection model for Galileo; (iii)
a 5G DL-TDOA estimator and uncertainty model that retain the correlation
introduced by a shared reference gNB; and (iv) a reproducible epoch-wise
assessment protocol separating positioning errors, model-conditional bounds,
and operational availability. These ingredients are established techniques;
scientific novelty must ultimately be demonstrated through the joint
comparison, modeling choices, and validated findings, rather than by claiming
a new protection-level formula.

\section{Integrity definitions and risk accounting}\label{sec:integrity}

Let the receiver's true position at epoch $k$ be $\vect{p}_k$ and the
estimated position be $\hat{\vect{p}}_k$, both expressed in local ENU
coordinates. Define the horizontal and vertical errors by
\begin{align}
 e_{H,k} &= \sqrt{(\hat E_k-E_k)^2+(\hat N_k-N_k)^2},\label{eq:eh}\\
 e_{V,k} &= |\hat U_k-U_k|.\label{eq:ev}
\end{align}

For protection levels $HPL_k$ and $VPL_k$, the target conditional
properties under an explicitly specified nominal model $\mathcal{M}_0$ are
\begin{align}
 \prob\!\left(e_{H,k}>HPL_k\mid\mathcal{M}_0\right)&\leq\alpha_H,\label{eq:h-risk}\\
 \prob\!\left(e_{V,k}>VPL_k\mid\mathcal{M}_0\right)&\leq\alpha_V.\label{eq:v-risk}
\end{align}
These inequalities do not automatically hold under unknown faults, NLOS
biases, model mismatch, or an incorrect solution branch. A total integrity
risk definition must also account for distinct fault hypotheses
$\mathcal{H}_j$, their occurrence probabilities, missed-detection events,
warning latency, and the operation's time horizon. Symbolically,
\begin{equation}\label{eq:total-risk}
 P_{\mathrm{HMI}}=\sum_j \prob(\mathcal{H}_j)\,
 \prob\!\left(\mathrm{hazardous\ error\ and\ no\ timely\ alert}\mid
 \mathcal{H}_j\right),
\end{equation}
where the hypotheses form a partition of the considered operating space.
Equation~\eqref{eq:total-risk} is a risk-accounting identity, not a
validated estimate of $P_{\mathrm{HMI}}$. For nonexclusive events,
appropriate union bounds are required.

When application-specific horizontal and vertical alert limits $HAL$
and $VAL$ are available, the corresponding epoch-wise operational
availability indicator is
\begin{equation}\label{eq:avail}
 a_k=\mathbb{1}\!\left\{\text{valid}_k,\ HPL_k\leq HAL,
 \ VPL_k\leq VAL\right\}.
\end{equation}
Without specified alert limits, only positioning/PL computability can be
reported; operational integrity availability is undefined.

\section{Simulation environment and reproducibility}\label{sec:scenario}

The framework reuses the urban pedestrian scenario developed in
\cite{SanchezHerederoThesis}. Building geometry and the walking trajectory
along the Plaza del Rey--Calle del Barquillo area of Madrid originate from
the existing simulation. A common local ENU frame is used for transmitter
geometry, estimated positions, and error evaluation. The original trajectory
is sampled at 2\,Hz and the Galileo positioning stage at 1\,Hz. The 5G
post-processing uses the epoch interval stored by its positioning simulation,
$T_{\mathrm{5G}}$, and aligns the corresponding truth and visibility arrays.
The exact number of evaluable epochs is determined from the exported files,
not inferred from the nominal path duration.

The 5G signal-generation stage is a representative standardized FR2
configuration rather than a reproduction of a commercial 5G deployment in
Madrid. The simulated transmitter locations are based on an existing
single-operator base-station geometry; the transmitted PRS waveform,
correlation-based timing measurements, and propagation simulation are
performed offline. Relevant radio parameters are listed in
Table~\ref{tab:scenario}; PRS signal structure and positioning procedures
follow \cite{ETSI38211,ETSI38305}, with urban macro propagation informed by
\cite{ETSI38901}.

\begin{table}[t]
\caption{Configuration inherited from the positioning simulation. These
are simulation parameters, not measured operational-network properties.}
\label{tab:scenario}
\begin{tabular}{@{}p{0.47\linewidth}p{0.46\linewidth}@{}}
\toprule
Parameter & Model configuration\\
\midrule
Coordinate frame & Local ENU\\
Trajectory / Galileo rate & 2\,Hz / 1\,Hz\\
5G positioning rate & $1/T_{\mathrm{5G}}$ from exported data\\
5G band / carrier frequency & FR2 n260 / 38.5\,GHz\\
NR subcarrier spacing & 120\,kHz ($\mu=3$)\\
5G channel bandwidth & 400\,MHz (264 PRBs)\\
PRS generation & 12 OFDM symbols, comb size 6\\
Propagation model & 3GPP urban macro, LOS selection\\
Galileo error model & Single frequency; see Sec.~\ref{sec:galileo}\\
\bottomrule
\end{tabular}
\end{table}

The original 5G positioning implementation exports
\texttt{epochEstim}, \texttt{tdoaEstim}, \texttt{tdoaError},
\texttt{LoSMatrix}, \texttt{LoSbStationsENU}, and the receiver ground truth.
The integrity script reconstructs an observation vector and transmitter
subset at each epoch. Importantly, \texttt{tdoaError} and ground-truth
positions are used only to validate the estimator \emph{after} computation;
they are not inputs to the operational protection-level calculation.

\section{Galileo positioning and integrity model}\label{sec:galileo}

\subsection{Pseudorange observation and least-squares solution}

For $n_k$ visible satellites at epoch $k$, let $\vect{s}_{i,k}$ denote
the position of satellite $i$, corrected to a common coordinate epoch, and
$\vect{p}_k$ the receiver position. The corrected pseudorange model is
\begin{equation}\label{eq:gnss-obs}
 \tilde\rho_{i,k}=\norm{\vect{p}_k-\vect{s}_{i,k}}+b_k
 +\varepsilon_{i,k}+f_{i,k},
\end{equation}
where $b_k=c\delta t_k$ is receiver clock bias in meters,
$\varepsilon_{i,k}$ collects nominal residual measurement errors, and
$f_{i,k}$ represents a measurement fault or unmodeled systematic bias.
Corrections for satellite clocks, relativistic effects, Earth rotation,
and atmosphere are assumed to precede estimation. This decomposition is
consistent with the original pseudorange model in
\cite{SanchezHerederoThesis}.

The estimated state in the local navigation frame is
$\vect{\theta}_k=[E_k,N_k,U_k,b_k]\trans$. If
$\vect{u}_{i,k}=(\vect{s}_{i,k}-\hat{\vect{p}}_k)/
\norm{\vect{s}_{i,k}-\hat{\vect{p}}_k}$ is the line-of-sight vector
from the user toward satellite $i$, each row of the geometry matrix is
\begin{equation}\label{eq:Ggnss}
 \mat{G}_{G,k}(i,:)=[-u_{E,i,k},-u_{N,i,k},-u_{U,i,k},1],
 \qquad \mat{G}_{G,k}\in\mathbb{R}^{n_k\times4}.
\end{equation}

After first-order linearization, $\Delta\vect{\rho}_k\simeq
\mat{G}_{G,k}\Delta\vect{\theta}_k+\vect{\varepsilon}_k$.
The baseline integrity implementation assumes mutually independent,
identically distributed Gaussian pseudorange errors with
$\sigma_G=4.6\,\mathrm{m}$, an approximate LOS uncertainty adopted in the
original thesis, informed by the Galileo Open Service error budget\cite{GalileoSDD}. This is a \emph{simulation assumption}, not an
environment-independent bound or an assertion that urban multipath is
Gaussian. Under this assumption,
\begin{align}
 \mat{A}_{G,k}&=(\mat{G}_{G,k}\trans\mat{G}_{G,k})^{-1}
 \mat{G}_{G,k}\trans,\label{eq:Agnss}\\
 \mat{P}_{G,k}&=\sigma_G^2
 (\mat{G}_{G,k}\trans\mat{G}_{G,k})^{-1}.\label{eq:Pgnss}
\end{align}
A valid three-dimensional position and clock estimate requires
$n_k\geq4$ and $\operatorname{rank}(\mat{G}_{G,k})=4$.
In computation, QR or singular-value decomposition should replace
explicit inversion when matrices are poorly conditioned.

\subsection{Nominal Gaussian horizontal and vertical bounds}

Let $\mat{P}_{G,k}^{EN}$ be the East--North block of the position
covariance, and $P_{G,k}^{UU}$ its Up variance. Define the two-sided
vertical Gaussian quantile $q_V(\alpha)=\Phi^{-1}(1-\alpha/2)$, where
$\Phi$ denotes the standard-normal cumulative distribution function.
A conservative two-dimensional radial bound follows from the largest
horizontal covariance eigenvalue:
\begin{align}
 HPL^{(0)}_{G,k}&=
 \sqrt{-2\log\alpha_H}\,
 \sqrt{\lambda_{\max}(\mat{P}_{G,k}^{EN})},\label{eq:hplgnssnom}\\
 VPL^{(0)}_{G,k}&=
 q_V(\alpha_V)\sqrt{P_{G,k}^{UU}}.\label{eq:vplgnssnom}
\end{align}
These quantities are \emph{fault-free, local, model-conditional} bounds.
They do not by themselves provide Galileo RAIM protection against faulty
measurements.

\subsection{Residual-based RAIM fault detection}\label{sec:raim}

Define the orthogonal projection onto the modeled measurement subspace
and its residual complement:
\begin{equation}\label{eq:gnss-BQ}
 \mat{B}_k=\mat{G}_{G,k}\mat{A}_{G,k},\qquad
 \mat{Q}_k=\mat{I}_{n_k}-\mat{B}_k.
\end{equation}
With post-fit pseudorange residuals
$\vect{v}_k=\tilde{\vect{\rho}}_k-
\vect{h}(\hat{\vect{\theta}}_k)$,
the test statistic is
\begin{equation}\label{eq:raim-stat}
 T_k=\frac{\vect{v}_k\trans\vect{v}_k}{\sigma_G^2}.
\end{equation}
Under the fault-free Gaussian linearized model,
\begin{equation}\label{eq:raim-null}
 T_k\sim\chi^2_{\nu_k},\qquad \nu_k=n_k-4.
\end{equation}
For $n_k\leq4$, the residual degrees of freedom are nonpositive or zero,
so no informative residual consistency test is available. With at least
five linearly independent satellite observations and
$P_{FA}\in(0,1)$, the detection threshold is
\begin{equation}\label{eq:raim-thresh}
 \gamma_k=F^{-1}_{\chi^2_{\nu_k}}(1-P_{FA}),
 \qquad \text{alarm if }T_k>\gamma_k.
\end{equation}

For a single range bias of amplitude $a_i$ on satellite $i$,
$\vect{f}=a_i\vect{e}_i$, the test statistic follows a noncentral
chi-square distribution:
\begin{align}
 T_k\mid(\vect{f}=a_i\vect{e}_i)
 &\sim\chi'^2_{\nu_k}(\lambda_{i,k}),\label{eq:raim-alt}\\
 \lambda_{i,k}&=\frac{a_i^2}{\sigma_G^2}\,Q_{k,ii},
 \label{eq:raim-lambda}\\
 P_{MD,i,k}(a_i)&=F_{\chi'^2_{\nu_k}(\lambda_{i,k})}
 (\gamma_k).\label{eq:raim-pmd}
\end{align}
The last equation is the missed-detection probability of the specified
single-bias hypothesis, conditional on the stated model and threshold.

\subsection{Fault sensitivity and limitations of a slope-based bound}

The linearized position effect of a single-satellite bias is
$\Delta\vect{\theta}_{f,k}\simeq\mat{A}_{G,k}\vect{f}$.
For $Q_{k,ii}>0$ define horizontal and vertical slopes
\begin{align}
 S_{H,i,k}&=\frac{\sqrt{A_{G,k}(1,i)^2+A_{G,k}(2,i)^2}}
 {\sqrt{Q_{k,ii}}},\label{eq:slopeH}\\
 S_{V,i,k}&=\frac{|A_{G,k}(3,i)|}
 {\sqrt{Q_{k,ii}}}.\label{eq:slopeV}
\end{align}
A noncentrality parameter $\lambda_\star$ can be selected such that
\begin{equation}\label{eq:lambda-star}
 F_{\chi'^2_{\nu_k}(\lambda_\star)}(\gamma_k)=P_{MD}^{\star},
\end{equation}
where $P_{MD}^{\star}$ is an \emph{allocated} missed-detection
probability. The corresponding detection-sensitivity quantities are
\begin{align}
 \Delta_{H,k}^{\rm MDE}&=\sigma_G\sqrt{\lambda_\star}
 \max_i S_{H,i,k},\label{eq:raim-Hmde}\\
 \Delta_{V,k}^{\rm MDE}&=\sigma_G\sqrt{\lambda_\star}
 \max_i S_{V,i,k}.\label{eq:raim-Vmde}
\end{align}
These are maximum position biases at the \emph{minimum detectable error}
(MDE) condition, not automatically valid fault protection levels. As an
intermediate engineering diagnostic, the quantities
\begin{equation}\label{eq:raim-proxy}
 \widetilde{HPL}_{G,k}=HPL^{(0)}_{G,k}+\Delta_{H,k}^{\rm MDE},\quad
 \widetilde{VPL}_{G,k}=VPL^{(0)}_{G,k}+\Delta_{V,k}^{\rm MDE}
\end{equation}
may be plotted, but must explicitly be labeled \emph{RAIM sensitivity
proxies}. Their interpretation as formal integrity protection levels
would additionally require a complete fault-frequency model, an
appropriate joint risk bound for undetected errors and the residual test,
and verification of the adopted risk allocations
\cite{Brown1992,WalterEnge1995}.

If independent physical bounds $|f_i|\leq\bar f_i$ are justified,
an alternative, conditional deterministic-bias bound is
\begin{equation}\label{eq:gnss-bounded}
 HPL_{G,k}^{\mathrm{bounded}}
 =HPL_{G,k}^{(0)}+
 \sum_i \bar f_i\norm{\mat{A}_{G,k}(1\!:\!2,i)}_2,
\end{equation}
with the analogous vertical bound
$VPL_{G,k}^{\mathrm{bounded}}=VPL_{G,k}^{(0)}+
\sum_i \bar f_i|A_{G,k}(3,i)|$.
Unlike the MDE proxy, this bound uses explicit limits on all modeled
simultaneous additive biases, but it remains local and invalid for
unbounded faults.

\section{5G NR DL-TDOA integrity model}\label{sec:5g}

\subsection{Timing observation model and reference differencing}

Consider $M_k$ visible terrestrial transmission--reception points (TRPs)
with known ENU coordinates $\vect{s}_{i,k}$. Let
$r_{i,k}(\vect{p})=\norm{\vect{p}-\vect{s}_{i,k}}$.
The following timing model presumes that known PRS transmission
scheduling offsets have been removed by the signal-processing chain.
This precondition must be checked against the waveform implementation;
otherwise a scheduling residual constitutes an additional bias.
The estimated TOA is modeled in seconds as
\begin{equation}\label{eq:toa5g}
 \tilde\tau_{i,k}=\frac{r_{i,k}(\vect{p}_k)}{c}
 +b_{UE,k}+b_{gNB,i,k}+\varepsilon_{i,k}+\beta_{i,k},
\end{equation}
where $b_{UE,k}$ is common UE clock error, $b_{gNB,i,k}$ is transmitter
timing error relative to the adopted network timescale,
$\varepsilon_{i,k}$ is random timing error, and $\beta_{i,k}$ covers
additional deterministic or model-bounded propagation bias. The common
receiver clock cancels in a difference, provided it is the same across
measurements within the considered observation interval; relative
gNB timing biases do not cancel.

Taking gNB 1 as reference, define the $(M_k-1)\times M_k$ differencing matrix
\begin{equation}\label{eq:D}
 \mat{D}_k=\begin{bmatrix}-\vect{1}_{M_k-1}&\mat{I}_{M_k-1}\end{bmatrix},
\end{equation}
and the measured range-difference vector
\begin{equation}\label{eq:ztdoa}
 \vect{z}_k=c\mat{D}_k\tilde{\vect{\tau}}_k.
\end{equation}
The ideal measurement function is
\begin{equation}\label{eq:htdoa}
 h_{i1,k}(\vect{p})=r_{i,k}(\vect{p})-r_{1,k}(\vect{p}),
 \qquad i=2,\ldots,M_k.
\end{equation}
This sign convention matches the exported simulation values
$\mathrm{TDOA}_{i1}=\mathrm{TOA}_i-\mathrm{TOA}_1$.

\subsection{Geometry matrix and correlated observation covariance}

Linearizing \eqref{eq:htdoa} around $\hat{\vect{p}}_k$ gives
\begin{equation}\label{eq:linear5g}
 \vect{z}_k-\vect{h}_k(\hat{\vect{p}}_k)
 \simeq\mat{G}_{5,k}\,\Delta\vect{p}_k+\vect{\eta}_k,
\end{equation}
with rows
\begin{equation}\label{eq:G5g}
 \mat{G}_{5,k}(i-1,:)=
 \frac{(\hat{\vect{p}}_k-\vect{s}_{i,k})\trans}
 {r_{i,k}(\hat{\vect{p}}_k)}-
 \frac{(\hat{\vect{p}}_k-\vect{s}_{1,k})\trans}
 {r_{1,k}(\hat{\vect{p}}_k)}.
\end{equation}
At least four gNBs and
$\operatorname{rank}(\mat{G}_{5,k})=3$ are needed for an
instantaneous unconstrained three-dimensional solution. Four gNBs
provide no redundant TDOA measurements for residual-based fault detection.

Let $\mat{C}_{\tau,k}$ denote the covariance of the individual timing
errors, in seconds squared. Their differencing produces
\begin{equation}\label{eq:R5g}
 \boxed{\mat{R}_{5,k}=c^2\mat{D}_k\mat{C}_{\tau,k}
 \mat{D}_k\trans}.
\end{equation}
For uncorrelated TOAs with a common standard deviation
$\sigma_\tau$, the result simplifies to
\begin{equation}\label{eq:R5gsimple}
 \mat{R}_{5,k}=c^2\sigma_\tau^2
 \left(\mat{I}_{M_k-1}+\vect{1}_{M_k-1}\vect{1}_{M_k-1}\trans\right).
\end{equation}
Thus diagonal elements equal $2c^2\sigma_\tau^2$ and distinct off-diagonal
elements equal $c^2\sigma_\tau^2$. The nonzero cross-covariances arise
from reuse of the reference TOA. Treating all TDOA measurements as
independent would generally lead to an incorrect position covariance.

\subsection{Nonlinear weighted least squares and position covariance}

The 5G position is estimated by minimizing
\begin{equation}\label{eq:wls5g}
 \hat{\vect{p}}_k=\arg\min_{\vect{p}}
 (\vect{z}_k-\vect{h}_k(\vect{p}))\trans
 \mat{R}_{5,k}^{-1}
 (\vect{z}_k-\vect{h}_k(\vect{p})).
\end{equation}
A Gauss--Newton step, evaluated at iterate $\vect{p}^{(t)}$, is
\begin{equation}\label{eq:gnstep}
 \Delta\vect{p}^{(t)}=
 \left(\mat{G}^{\mathsf T}\mat{R}^{-1}\mat{G}\right)^{-1}
 \mat{G}\trans\mat{R}^{-1}
 \left(\vect{z}-\vect{h}(\vect{p}^{(t)})\right).
\end{equation}
The actual numerical implementation uses Cholesky whitening, QR
factorization, and a backtracking line search rather than explicit
matrix inversion. After convergence, the local position covariance
and measurement-to-position sensitivity matrix are
\begin{align}
 \mat{P}_{5,k}&\simeq
 (\mat{G}_{5,k}\trans\mat{R}_{5,k}^{-1}\mat{G}_{5,k})^{-1},
 \label{eq:P5g}\\
 \mat{K}_{5,k}&=
 (\mat{G}_{5,k}\trans\mat{R}_{5,k}^{-1}\mat{G}_{5,k})^{-1}
 \mat{G}_{5,k}\trans\mat{R}_{5,k}^{-1}.\label{eq:K5g}
\end{align}
The covariance approximation assumes convergence to the correct solution
branch, local linearity, correct $\mat{R}_{5,k}$, and no unmodeled bias.
Very poor geometry, rank deficiency, or failed numerical convergence
causes the epoch to be flagged invalid.

\subsection{Gaussian and bounded-bias protection levels}

Let $\vect{\beta}_{\tau,k}$ denote the total uncorrected timing bias,
including residual gNB synchronization error and modeled NLOS excess
delay, in seconds. The propagated local position error is
\begin{equation}\label{eq:Ktoa}
 \Delta\vect{p}_{\beta,k}\simeq
 \underbrace{c\mat{K}_{5,k}\mat{D}_k}_{\mat{K}_{\tau,k}}
 \vect{\beta}_{\tau,k}.
\end{equation}
Assume justified componentwise bounds
$|\beta_{\tau,j,k}|\leq \bar\beta_{j,k}$ and define
$\mat{K}_{\tau,k}^{EN}$ as the first two rows of $\mat{K}_{\tau,k}$.
The triangle inequality yields conservative bias contributions
\begin{align}
 B_{H,k}&=\sum_{j=1}^{M_k}\bar\beta_{j,k}
 \norm{\mat{K}_{\tau,k}^{EN}(:,j)}_2,\label{eq:BH5g}\\
 B_{V,k}&=\sum_{j=1}^{M_k}\bar\beta_{j,k}
 |K_{\tau,k}(3,j)|.\label{eq:BV5g}
\end{align}
For the Gaussian part of the position error, the horizontal radial
and vertical two-sided bounds become
\begin{align}
 HPL_{5,k}&=B_{H,k}+
 \sqrt{-2\log\alpha_H}\,
 \sqrt{\lambda_{\max}(\mat{P}_{5,k}^{EN})},
 \label{eq:HPL5g}\\
 VPL_{5,k}&=B_{V,k}+
 \Phi^{-1}(1-\alpha_V/2)\sqrt{P_{5,k}^{UU}}.
 \label{eq:VPL5g}
\end{align}
For bounded deterministic bias and zero-mean Gaussian random error,
these are conservative, \emph{local model-conditional} bounds.
A mathematical derivation is provided in Appendix~\ref{app:bounds}.
An unbounded NLOS error, a clock jump outside the assumed bound, or a
covariance mismatch invalidates the stated probability coverage.

In the current implementation, the illustrative parameters are
$\sigma_\tau=1\,\mathrm{ns}$ and
$\alpha_H=\alpha_V=5\times10^{-8}$; the synchronization and
NLOS bias bounds are initialized to zero. These are explicitly
\emph{uncalibrated scenario inputs}. A zero bias bound represents a
fault-free hypothesis, not evidence that synchronization or NLOS
errors vanish. The two tail allocations sum to $10^{-7}$ for the
combined nominal horizontal/vertical exceedance event using a union
bound; they do not establish an unconditional integrity risk of
$10^{-7}$ per epoch or per mission.

\section{Epoch-wise implementation and validation protocol}\label{sec:method}

The proposed assessment is implemented as post-processing of the original
Galileo and 5G simulation outputs; this separation avoids changing the
underlying signal-generation experiments. At each epoch, the 5G routine
uses the gNBs marked visible in \texttt{LoSMatrix}; the first visible gNB
is the reference corresponding to the stored \texttt{tdoaEstim} row.
The reference gNB may change between epochs. Because only TDOA
measurements are stored, a synthetic TOA vector
$[0,\Delta\tau_{21},\ldots,\Delta\tau_{M1}]\trans$ is used.
Its common time origin is arbitrary, but multiplying by $\mat{D}_k$
recovers the observed differences exactly. This reconstruction does
\emph{not} recover independent physical TOAs or calibrate their covariance.

A valid 5G positioning epoch requires a compatible observation subset,
adequate geometric rank, finite inputs, a positive-definite TDOA
covariance, and successful WLS convergence. The result includes
$\hat{\vect{p}}_k$, $\mat{G}_{5,k}$, $\mat{R}_{5,k}$,
$\mat{P}_{5,k}$, $HPL_{5,k}$, $VPL_{5,k}$,
the selected reference gNB, the number of measurements,
and a diagnostic status flag. The original 5G position estimate is
used as the \emph{initial guess} only; the reported protection levels
correspond to the newly computed WLS estimate, not the original
algebraic TDOA solution.

\begin{algorithm}[t]
\caption{Per-epoch 5G DL-TDOA conditional integrity computation}
\label{alg:5g}
\begin{algorithmic}[1]
\Require Saved TDOA data, gNB geometry, initial 5G positions,
$\mat{C}_{\tau,k}$, $\bar{\vect{\beta}}_k$, $\alpha_H$, $\alpha_V$
\For{each 5G positioning epoch $k$}
  \State Determine visible gNB indices and epoch-specific reference
  \If{fewer than four usable gNBs or inconsistent TDOA indices}
    \State Mark epoch invalid; \textbf{continue}
  \EndIf
  \State Reconstruct equivalent relative TOA vector
  \State Form $\mat{D}_k$, $\vect{z}_k$ and $\mat{R}_{5,k}$
  \State Solve nonlinear WLS using whitening and line search
  \If{geometry invalid or WLS does not converge}
    \State Mark epoch invalid; \textbf{continue}
  \EndIf
  \State Compute $\mat{G}_{5,k}$, $\mat{P}_{5,k}$ and $\mat{K}_{\tau,k}$
  \State Compute $HPL_{5,k}$ and $VPL_{5,k}$ using
    Eqs.~\eqref{eq:HPL5g}--\eqref{eq:VPL5g}
  \State Store validity, PLs, covariance, reference and geometry
  \State Separately compute truth-based errors for retrospective validation
\EndFor
\end{algorithmic}
\end{algorithm}

Galileo epochs are processed analogously, but with the Galileo
pseudorange geometry matrix, at least four satellites for LS, and
at least five for residual-based RAIM detection. Under the simplified
single-fault model, the residual test and MDE sensitivity analysis are
computed only when the requisite redundancy is available.

\subsection{Evaluation indicators}\label{sec:metrics}

Let $\mathcal{V}_S$ be the set of valid solution epochs for technology
$S\in\{G,5\}$, and let $N_{\mathrm{all},S}$ be the total number of
epochs considered. The following quantities will be reported:
\begin{align}
 A_{\mathrm{comp},S}&=\frac{|\mathcal{V}_S|}{N_{\mathrm{all},S}},
 \label{eq:Acomp}\\
 \widehat p_{H,S}&=\frac{1}{|\mathcal{V}_S|}
 \sum_{k\in\mathcal{V}_S}\mathbb{1}\{e_{H,k}>HPL_{S,k}\},
 \label{eq:pHobs}\\
 \widehat p_{V,S}&=\frac{1}{|\mathcal{V}_S|}
 \sum_{k\in\mathcal{V}_S}\mathbb{1}\{e_{V,k}>VPL_{S,k}\}.
 \label{eq:pVobs}
\end{align}
Observed exceedance frequencies are descriptive diagnostics and are not
proofs of target probabilities as small as $10^{-8}$.
The interpretation depends on the available number of independent
trials, temporal correlation, and uncertainty of the underlying
fault and measurement distributions.

Further indicators are the median and 95th percentile of valid-epoch
$HPL$ and $VPL$, the proportion of epochs satisfying specified
$HAL$/$VAL$, the number of gNBs or satellites used, and the
condition number of the geometry matrix. Results should be stratified
by LOS count, geometry quality, and error-model scenario. A residual
threshold alert is reported separately from a numerical PL
exceedance observed with ground truth.

\section{Extension to hybrid Galileo--5G filtering}\label{sec:ekf}

The master's thesis already implements a state-space EKF combining the
positioning modalities \cite{SanchezHerederoThesis}. The integrity-aware
EKF analysis is \emph{planned but not yet implemented or validated}.
One consistent state parameterization, with receiver clock bias in
meters, is
\begin{equation}\label{eq:ekf-state}
 \vect{x}_k=[E_k,N_k,U_k,v_{E,k},v_{N,k},v_{U,k},b_k]\trans.
\end{equation}
For a constant-velocity process model,
\begin{align}
 \hat{\vect{x}}_{k}^{-}&=\mat{F}_k\hat{\vect{x}}_{k-1}^{+},
 \label{eq:ekf-pred1}\\
 \mat{P}_{k}^{-}&=\mat{F}_k\mat{P}_{k-1}^{+}\mat{F}_k\trans
 +\mat{Q}_k.\label{eq:ekf-pred2}
\end{align}
A generic measurement update uses the innovation
$\vect{y}_k=\vect{z}_k-\vect{h}(\hat{\vect{x}}_k^-)$,
innovation covariance $\mat{S}_k=\mat{H}_k\mat{P}_k^-
\mat{H}_k\trans+\mat{R}_k$, and gain
$\mat{L}_k=\mat{P}_k^-\mat{H}_k\trans\mat{S}_k^{-1}$.
Numerically, the Joseph covariance update is preferred:
\begin{equation}\label{eq:joseph}
 \mat{P}_k^+=(\mat{I}-\mat{L}_k\mat{H}_k)\mat{P}_k^-
 (\mat{I}-\mat{L}_k\mat{H}_k)\trans+
 \mat{L}_k\mat{R}_k\mat{L}_k\trans.
\end{equation}
For Galileo pseudoranges, the clock-bias column of the Jacobian is
unity; for 5G range-difference measurements, that column is zero.
The 5G measurement-noise matrix must retain the common-reference
correlation in Eq.~\eqref{eq:R5g}.

\emph{Crucially}, applying the instantaneous Gaussian PL equations to
$\mat{P}_k^+$ alone does not establish filter integrity. An EKF can
propagate faults over time; temporal dependence, prior-state errors,
process-model mismatch, innovations, and fault-detection/exclusion
hypotheses must be modeled. These issues form a separate research
extension, not a demonstrated result of this paper.

\section{Results: predefined analyses awaiting simulation output}\label{sec:results}

This section defines the intended results and presentation format.
No quantitative results from the new integrity analysis are available
for inclusion in the current manuscript. All numerical cells and
figure panels below are deliberately left unpopulated.

\subsection{Positioning and integrity availability}

Table~\ref{tab:avail-results} will report valid-epoch fractions and
availability under explicit application alert limits. The distinction
between insufficient geometry, solver failure, and exceeded PL shall
be preserved.

\begin{table}[t]
\caption{Planned availability comparison. Numerical results pending.
$HAL$ and $VAL$ must be selected before computing operational availability.}
\label{tab:avail-results}
\begin{tabular}{@{}lccc@{}}
\toprule
Metric & Galileo LS & Galileo RAIM & 5G WLS\\
\midrule
Total evaluated epochs & TBD & TBD & TBD\\
Valid position epochs & TBD & TBD & TBD\\
Valid covariance/PL epochs & TBD & TBD & TBD\\
RAIM test-available epochs & N/A & TBD & N/A\\
PL below $HAL$/$VAL$ & TBD & TBD & TBD\\
Computability, \% & TBD & TBD & TBD\\
Operational availability, \% & TBD & TBD & TBD\\
\bottomrule
\end{tabular}
\end{table}

\subsection{Protection levels versus actual error}

Figures~\ref{fig:horizontal-placeholder} and
\ref{fig:vertical-placeholder} reserve space for the epoch-wise
horizontal and vertical position errors and their corresponding
model-conditional protection levels. For Galileo, the plotted
RAIM MDE sensitivity proxy must be visibly distinguished from the
nominal bound and from any later validated integrity PL.

\begin{figure}[t]
\centering
\fbox{\begin{minipage}[c][36mm][c]{0.91\linewidth}
\centering \textbf{Horizontal error and HPL versus epoch}\\[2mm]
\emph{Insert Galileo and 5G plots from the validated simulation.}\\
\emph{Optional: $HAL$ and fault-alarm markers.}
\end{minipage}}
\caption{Planned time-series comparison of horizontal positioning error and
model-conditional HPL for Galileo and 5G. Data and curves pending.}
\label{fig:horizontal-placeholder}
\end{figure}

\begin{figure}[t]
\centering
\fbox{\begin{minipage}[c][36mm][c]{0.91\linewidth}
\centering \textbf{Vertical error and VPL versus epoch}\\[2mm]
\emph{Insert Galileo and 5G plots from the validated simulation.}\\
\emph{Optional: $VAL$ and geometry-validity markers.}
\end{minipage}}
\caption{Planned time-series comparison of absolute vertical error and
model-conditional VPL for Galileo and 5G. Data and curves pending.}
\label{fig:vertical-placeholder}
\end{figure}

\begin{table}[t]
\caption{Planned protection-level coverage and accuracy metrics, conditioned
on valid-epoch solutions. No numerical results have been inserted.}
\label{tab:pl-results}
\begin{tabular}{@{}p{0.43\linewidth}p{0.24\linewidth}p{0.24\linewidth}@{}}
\toprule
Metric & Galileo & 5G DL-TDOA\\
\midrule
Horizontal RMSE (m) & TBD & TBD\\
Median / 95th-percentile HPL (m) & TBD / TBD & TBD / TBD\\
Median / 95th-percentile VPL (m) & TBD / TBD & TBD / TBD\\
Horizontal exceedances (count) & TBD & TBD\\
Vertical exceedances (count) & TBD & TBD\\
PL/geometry rejection epochs (count) & TBD & TBD\\
\bottomrule
\end{tabular}
\end{table}

\subsection{Geometry and failure sensitivity}

The analysis will study the relation between transmitter visibility,
geometry conditioning, and protection levels. Particular attention will
be given to epochs with exactly four gNBs, where 5G has no redundant
TDOA observations, and to Galileo epochs with four satellites, where
position may be available without residual RAIM detectability.

For Galileo, the residual statistic $T_k$, threshold $\gamma_k$,
$P_{FA}$, $P_{MD}^{\star}$, and worst-case slope will be evaluated under
explicit injected single-satellite bias scenarios. A separate study
should inject bounded gNB timing offsets and NLOS delays to quantify
how 5G bounds deteriorate as assumptions are relaxed.

\subsection{Impact of TDOA covariance modeling}

The simulation will compare three analyses using the same saved
observations: (i) the original algebraic 5G position estimate,
(ii) WLS with the full, common-reference covariance
$\mat{R}_{5,k}$, and (iii) a deliberately misspecified WLS model
that discards the off-diagonal covariance terms. Only the full-
covariance model is adopted for the main integrity results; the
diagonal approximation is included as a controlled ablation.
Table~\ref{tab:cov-ablation} defines the intended comparison.

\begin{table}[t]
\caption{Planned covariance-model ablation (all results pending).}
\label{tab:cov-ablation}
\begin{tabular}{@{}lccc@{}}
\toprule
Method & Horizontal RMSE & Median HPL & HPL exceedances\\
\midrule
Original 5G estimator & TBD & N/A & N/A\\
WLS, full $\mat{R}_{5,k}$ & TBD & TBD & TBD\\
WLS, diagonal $\mat{R}_{5,k}$ & TBD & TBD & TBD\\
\bottomrule
\end{tabular}
\end{table}

\subsection{Statistical evaluation plan}

The following experiment families are reserved for quantitative study:
(i) LOS-only Gaussian noise with calibrated variances; (ii) unequal
per-link variances and shared reference effects; (iii) bounded
synchronization errors; (iv) NLOS bias distributions or worst-case
bounds; and (v) single-fault Galileo RAIM detection sweeps. Sensitivity
plots will vary $\sigma_G$, $\sigma_\tau$, $\alpha_H$,
$\alpha_V$, and fault assumptions independently. Any estimate of
rare integrity risk will require a statistically justified evaluation
procedure; a short single-trajectory simulation is insufficient to
verify probabilities on the order of $10^{-8}$.

\section{Discussion}\label{sec:discussion}

The two sensing modalities share a least-squares foundation, but their
measurement transformations differ. Galileo pseudoranges explicitly
include the receiver clock state, whereas DL-TDOA removes the common
receiver clock bias at the cost of correlated difference errors and
sensitivity to network synchronization. For the same ENU frame,
comparison of the corresponding position covariance blocks is valid
only after measurement units and covariance conventions are matched.

A common scalar $\sigma_G$ simplifies Galileo simulation analysis, but
urban canyon conditions do not imply identically distributed or
zero-mean pseudorange errors. In particular, the thesis's elevated
NLOS variance approximation does not guarantee that excess-path
biases are bounded or symmetric. Similarly, applying a fixed
$\sigma_\tau=1\,\mathrm{ns}$ to every gNB is an illustrative
configuration, not an experimentally calibrated statement about
PRS delay-estimation uncertainty. These are major threats to
integrity validity, not merely implementation details.

A key distinction is required between observability and fault
monitorability. Four satellites (including receiver-clock unknown)
or four gNBs (three independent TDOAs) may suffice for a position,
but do not supply residual redundancy for a global consistency test.
Moreover, low post-fit residuals do not exclude coherent or
weakly observable errors. Rigorous operational integrity therefore
requires explicit fault hypotheses and alert procedures, consistent
with the broader integrity literature \cite{Brown1992,Bartoletti2025}.

The intended scientific value lies in empirically quantifying the
tradeoffs among geometry, timing uncertainty, fault detectability,
and protection-level availability in the same urban corridor.
Neither new mathematical equations nor an integrity improvement
relative to prior work can be asserted until those comparisons are
performed and independently verified.

\section{Conclusions and future work}\label{sec:conclusions}

An epoch-wise comparative integrity-assessment framework has been
formulated for Galileo least-squares pseudorange positioning and
5G NR DL-TDOA weighted least squares. The Galileo methodology
combines local covariance analysis with a residual-based single-fault
RAIM model. The 5G methodology propagates correlated common-reference
TDOA uncertainty and explicitly bounded timing biases into
horizontal and vertical conditional protection levels.

The current contribution is a specified and implemented analysis
workflow rather than a demonstrated integrity-performance advantage.
Future work will calibrate timing and pseudorange error models,
inject realistic multipath and synchronization faults, perform
risk-consistent RAIM evaluation, quantify actual model-bound
exceedances and availability, and extend the assessment to the
hybrid EKF with temporal fault propagation. Conclusions on
comparative performance will be added only after results are
reproduced and validated.

\backmatter

\bmhead{Acknowledgements}
\todo{Confirm acknowledgements, grants and institutional contributions.}

\section*{Declarations}
\textbf{Funding.} \todo{Declare funding or state none.}

\textbf{Conflict of interest.} \todo{Confirm.}

\textbf{Data availability.} The mathematical formulation and intended data
products are described herein. \todo{Insert archived dataset identifier or repository
URL and reproducibility conditions after validation.}

\textbf{Code availability.} A MATLAB positioning simulation and an
associated integrity post-processing implementation are used.
\todo{Insert stable repository commit or archival DOI before submission.}

\textbf{Author contributions.} \todo{Complete after authorship is finalized.}

\begin{appendices}
\section{Derivation of TDOA correlation and Gaussian position bounds}
\label{app:bounds}

For identical independent timing errors,
$\mat{C}_\tau=\sigma_\tau^2\mat{I}_M$,
Eq.~\eqref{eq:R5g} gives
\begin{align}
 \mat{D}\mat{D}\trans&=
 \mat{I}_{M-1}+\vect{1}_{M-1}\vect{1}_{M-1}\trans,\\
 \operatorname{Var}(c(\varepsilon_i-\varepsilon_1))&=
 2c^2\sigma_\tau^2,\\
 \operatorname{Cov}(c(\varepsilon_i-\varepsilon_1),
 c(\varepsilon_j-\varepsilon_1))&=
 c^2\sigma_\tau^2\quad(i\ne j).
\end{align}

Let the horizontal random positioning error be
$\vect{w}_{EN}\sim\mathcal{N}(\vect{0},\mat{P}^{EN})$.
With $\vect{z}\sim\mathcal{N}(\vect{0},\mat{I}_2)$ we may write
$\vect{w}_{EN}=(\mat{P}^{EN})^{1/2}\vect{z}$, so
\begin{equation}\label{eq:radial-proof}
 \norm{\vect{w}_{EN}}_2\leq
 \sqrt{\lambda_{\max}(\mat{P}^{EN})}\norm{\vect{z}}_2.
\end{equation}
Because $\norm{\vect{z}}_2^2\sim\chi^2_2$ and
$\prob(\chi_2^2>x)=\exp(-x/2)$, a sufficient radius for tail
probability $\alpha_H$ is
$\sqrt{\lambda_{\max}(\mat{P}^{EN})}\sqrt{-2\ln\alpha_H}$.
With bounded biases, $\norm{\vect{w}_{EN}+\vect{b}_{EN}}
\leq\norm{\vect{w}_{EN}}+\norm{\vect{b}_{EN}}$.
Applying the triangle inequality again to
$\vect{b}_{EN}=\sum_j\mat{K}_\tau^{EN}(:,j)\beta_j$
proves the addition of $B_H$ in Eq.~\eqref{eq:HPL5g}.
The analogous scalar Gaussian inequality proves
Eq.~\eqref{eq:VPL5g}. All bounds use a first-order local
measurement-to-position mapping.

\section{Residual RAIM statistic under a single range bias}
\label{app:raim}

For full column rank $\mat{G}_G$ with $n>4$, the matrix
$\mat{Q}=\mat{I}-\mat{G}_G(\mat{G}_G\trans\mat{G}_G)^{-1}
\mat{G}_G\trans$ is symmetric and idempotent, with
$\operatorname{rank}(\mat{Q})=n-4$. In the linearized model,
$\vect{v}\simeq\mat{Q}(\vect{\varepsilon}+a_i\vect{e}_i)$.
The whitened quadratic form
$\vect{v}\trans\vect{v}/\sigma_G^2$ is therefore central
chi-square under $a_i=0$ and noncentral chi-square with
\begin{equation}
 \lambda_i=\frac{a_i^2}{\sigma_G^2}\vect{e}_i\trans
 \mat{Q}\vect{e}_i=\frac{a_i^2 Q_{ii}}{\sigma_G^2}
\end{equation}
under a single additive fault. The local position error is
$\mat{A}_G\vect{\varepsilon}+a_i\mat{A}_G\vect{e}_i$,
which motivates the slope definitions in
Eqs.~\eqref{eq:slopeH} and~\eqref{eq:slopeV}. Fault detection
and the joint hazard probability remain distinct questions.

\end{appendices}

\begin{thebibliography}{99}
\bibitem{SanchezHerederoThesis}
D. S\'anchez-Heredero Mart\'inez,
\emph{5G Assisted Precise Positioning in GNSS},
Master's thesis manuscript, Master's in Wireless Communications,
2026. Underlying simulation supplied by the author.

\bibitem{GalileoSDD}
European Union Agency for the Space Programme,
\emph{Galileo Open Service: Service Definition Document},
Issue 1.3, November 2023.
\url{https://www.gsc-europa.eu/electronic-library/programme-reference-documents/galileo-in-force/open-service}.

\bibitem{Brown1992}
R. G. Brown, A baseline GPS RAIM scheme and a note on the equivalence
of three RAIM methods, \emph{NAVIGATION}, vol. 39, no. 3,
pp. 301--316, 1992.
\url{https://doi.org/10.1002/j.2161-4296.1992.tb02278.x}.

\bibitem{WalterEnge1995}
T. Walter and P. Enge, Weighted RAIM for precision approach,
\emph{Proc. 8th International Technical Meeting of the Satellite
Division of The Institute of Navigation (ION GPS 1995)},
pp. 1995--2004, September 1995.
\url{https://www.ion.org/publications/abstract.cfm?articleID=2524}.

\bibitem{ETSI38305}
ETSI, \emph{5G; NG Radio Access Network (NG-RAN); Stage 2
Functional Specification of User Equipment (UE) Positioning in NG-RAN},
ETSI TS 138 305 V18.7.0 (3GPP TS 38.305, Release 18), February 2026.
\url{https://www.etsi.org/deliver/etsi_ts/138300_138399/138305/18.07.00_60/ts_138305v180700p.pdf}.

\bibitem{ETSI38211}
ETSI, \emph{5G; NR; Physical Channels and Modulation},
ETSI TS 138 211 V18.7.0 (3GPP TS 38.211, Release 18), July 2025.
\url{https://www.etsi.org/deliver/etsi_ts/138200_138299/138211/18.07.00_60/ts_138211v180700p.pdf}.

\bibitem{ETSI38901}
ETSI, \emph{5G; Study on Channel Model for Frequencies from
0.5 to 100\,GHz}, ETSI TR 138 901 V18.1.0
(3GPP TR 38.901, Release 18), February 2026.
\url{https://www.etsi.org/deliver/etsi_tr/138900_138999/138901/18.01.00_60/tr_138901v180100p.pdf}.

\bibitem{Bartoletti2025}
S. Bartoletti, I. Palam\`a, L. Chiaraviglio, S. Modarres Razavi,
Y. Zhao, G. Bianchi, and N. Blefari Melazzi,
Positioning integrity via uncertainty quantification,
\emph{IEEE Transactions on Vehicular Technology},
vol. 74, no. 1, pp. 1224--1237, January 2025.
\url{https://doi.org/10.1109/TVT.2024.3457861}.
\end{thebibliography}

\end{document}
